\manualchapter{Control reference}{sec:controls}
Panel map: \textbf{1} pitch and pitch/FM inputs;
\textbf{2} FM depth;
\textbf{3} level and level CV;
\textbf{4} shape buttons and indicators;
\textbf{5} audio outputs.

Four matching columns run from tuning at the top through FM, level and
shape to the bottom outputs. The CV sockets sit alongside their controls. A column's shape button advances through Sine,
Triangle, NES Triangle, Sample+Hold, LFSR Long and LFSR Short. The context
menu selects a named shape directly.

\begin{tabularx}{\textwidth}{@{}l l X@{}}
\toprule Control & Range; default & Operation \\\midrule
Frequency & $\pm2.5$ oct.; 0 & C4 reference; V/OCT adds one octave per volt. \\
FM & $-1$--$+1$; 0 & Bipolar depth $a$ adds $aV/5$ octaves. \\
Level & 0--255; 255 & Multiplied by $V/10$, rounded and clipped to 8 bits. \\
Shape & Six choices & Button advances; context menu names each choice. \\
\bottomrule
\end{tabularx}

Pitch, FM and Level normal forward across the columns. Their first inputs
normal to 0, 5 and 10\,V respectively. With FM unpatched, its knob is a
fine-tuning control spanning an octave either side of the main setting.
At zero FM depth an external FM signal has no effect. Level is a unipolar
amplitude control; negative CV closes it rather than inverting the waveform.

\begin{table}[ht]\centering
\begin{tabularx}{\textwidth}{@{}lX@{}}
\toprule Shape & Use \\\midrule
Sine & A simple tonal reference for tuning or modulation. \\
Triangle & A brighter tonal source with a triangular shape. \\
NES Triangle & A stepped triangle for coarse digital bass. \\
Sample+Hold & Clocked random levels; tune the update rate. \\
LFSR Long / Short & Shift-register patterns with different repetition lengths. \\
\bottomrule
\end{tabularx}
\caption{Selectable oscillator shapes.}\label{tab:oscillator-shapes}
\end{table}
